\section{Introduction}
When directional demand varies, a fixed green split can leave a queue on one axis while service on the other goes unused. Adaptive timing can respond to this imbalance, but a controller demonstration must make two different events observable: selecting a timing plan and applying it. A credible report must identify the data and decision path behind a recommendation, then establish which plan actually governed the next cycle. Traffic engineering guidance \cite{koonce}, a Hong Kong trial \cite{hk}, and a city-scale adaptive timing analysis \cite{wu} motivate the demand problem.

The Traffic Timing Assistant targets straight-through movements on the north--south (NS) and east--west (EW) axes of one intersection. Three questions govern the assessment:
\begin{enumerate}
\item How does demand-aware selection affect queues and waiting relative to a fixed balanced plan? The available paired host summaries quantify aggregate queue burden; per-vehicle waiting records are pending.
\item Can an operator explain a recommendation? The design records the input values and traversed predicates or fallback reason. Logged decision paths are needed to audit fidelity.
\item Did the plan take effect at the intended boundary? The proposal reports a board application test. Time-aligned state observations are needed to verify boundary behavior and mismatches.
\end{enumerate}

The contribution is the design of an auditable path from traffic observations to constrained plan selection and verified phase execution. It combines a three-class selector and a fallback policy, a cycle-boundary handoff between ARM software and FPGA logic, and a common record that relates the recommendation to device state. The proposal reports working model and controller cores. Structured task intake, cross-source validation and unified reporting remain integration work.
